%% HYPOTHESIS / EXPLORATORY DOCUMENT\n%% NOTE: Auto-generated exploratory hypotheses. NOT A PROOF.\n\documentclass{article}\n\usepackage{amsmath,amssymb}\n\begin{document}\n\title{Generated hypothesis from extracaffe.pdf (Hypothesis)}\maketitle\n\section{Context}\nPseudo-entropy: 0.534260\n\begin{verbatim}\n空間概念の量子化
Masaaki Yamaguchi

あらまし
始まりは、無限の時間の無が、場の理論として存在していた。この無が、オイラーの定数の大域的積分多様
体が、基本群による、相対性から、オイラーの定数の多様体積分としてのガンマ関数に化けて、ゼウスとなり、
このゼウスから、質量差エネルギーとして、母なる Jones 多項式ができた。この Jones 多項式から、各種の生
成物として、宇宙と異次元が、D-brane を柱として、生み出された。この Jones 多項式から、中間子を媒介に
して、素粒子が出来て、世界が切り開かれた。この世界は、Zeta 関数を禅に対して、Beta 関数が終わりに来
るように、時間が２通りの流れとして、存在された。この時間の流れの中で、互いの相対性が消えて、空間概
念の量子化としての、すべての共時性が最終目標にされた。
登場人物は、
大域的積分多様体のガンマ関数
時間と空間

父

トキ子おばあちゃん

Jones 多項式 母と彩さん
D-brane

山中崇史くん

Zeta 関数 無
推移律の積分

先生たち

アカシックレコードのヒント
著者

姉たち

私

1

Entropy on manifold and differential of equation
Masaaki Yamaguchi

1 Global Differential Equation
Global area in integrate and differential of equation emerge with quantum of image function built
on. Thurston conjugate of theorem come with these equation, and heat equation of entropy developed
it’ equation with. These equation of entropy is common function on zeta function, eight of differential
structure is with fourth power, then gravity, strong boson, weak boson, and Maxwell theorem is common
of entropy equation. Euler-Lagrange equation is same entropy with quantum of material equation. This
spectrum focus is, eight of differential structure bind of two perceptive power. These more say, zeta
function resolved is gravity and antigravity on non-catastrophe on D-brane emerge with same power of
level in fourth of power. This quantum group of equation inspect with universe of space on imaginary
and reality of Laplace equation.
It is possibility of function for these equation to emerge with entropy of fluctuation in equation on
prime parammeter. These equation is built with non-commutative formula. And these equation is
emerged with fundamental formula. These equation group resolved by Schrödinger, Heisenberg equation
and D-brane. More spectrum focus is, resolved equation by zeta function bind with gravity, antigravity
and Maxwell theorem. More resolved is strong boson and weak boson to unite with. These two of power
integrate with three dimension on zeta function.
Imagi\end{verbatim}\n\section{Extracted equations}\n\begin{equation}\label{eq:ex0}\n. \bigoplus \nabla C^{+}_{-} \cong M_3, R \supset Q, R \cap Q, R \subset M_3, C^{+} \bigoplus M_n, E^{+} \cap R^{+}\n\end{equation}\n\begin{equation}\label{eq:ex1}\n\n\end{equation}\n\begin{equation}\label{eq:ex2}\n. {\nabla \over \Delta} \int x f(x) dx, {\nabla R \over \Delta f}, \square = 2{\int {(R + \nabla_{i} \nabla_{j} f)^2 \over -(R + \Delta f)}}e^{-f}dV \square = {\nabla R \over \Delta f}, {d \over dt}g_{ij} = \square \to {\nabla f \over \Delta x}, (R + |\nabla f|^2)dm \to -2(R + \nabla_{i} \nabla_{j} f)^2 e^{-f}dV x^n + y^n = z^n \to \nabla \psi^2 = 8 \pi G T^{\mu\nu}, f(x + y) \ge f(x) \circ f(y) \mathrm{im}f / \mathrm{ker}f = \partial f, \mathrm{ker}f = \partial f, \mathrm{ker}f / \mathrm{im}f \cong \partial f, \mathrm{ker}f = f^{-1}(x)xf(x) f^{-1}(x)xf(x) = \int \partial f(x) d(\mathrm{kerf}) \to \nabla f = 2 _{n}C_{r} = {}_{n}C_{n-r} \to \mathrm{im}f / \mathrm{ker}f \cong \mathrm{ker}f / \mathrm{im}f\n\end{equation}\n\begin{equation}\label{eq:ex3}\n. this equation\n\end{equation}\n\begin{equation}\label{eq:ex4}\n. V/W = R/C \sum^{\infty}_{k=0}{x^k \over a_k f^k}, W/V = C/R \sum^{\infty}_{k=0}{a_k f^k \over x^k} V/W \cong W/V \cong R/C(\sum^{\infty}_{r=0} {}_nC_{r})^{-1} \sum^{\infty}_{k=0} x^k This equation is diffrential equation, then\n\end{equation}\n\begin{equation}\label{eq:ex5}\nis included with\n\end{equation}\n\begin{equation}\label{eq:ex6}\nW/V = xF(x), \chi(x) = (-1)^k a_k, \\Gamma(x) = \int e^{-x} x^{1 -t}dx, \sum^{n}_{k=0}a_k f^k = (f^k)’ \sum^{n}_{k=0}a_k f^k = \sum^{\infty}_{k~0} {}_{n}C_{r} f^k = (f^k)’, \sum^{\infty}_{k=0} a_k f^k = [f(x)], \sum^{\infty}_{k=0} a_k f^k = \alpha, \sum^{\infty}_{k=0} {1 \over a_k f^k}, \sum^{\infty}_{k=0} (a_k f^k)^{-1} = {1 \over 1 - z} {\int \int {1 \over (x \log x)(y \log y)}dxy} = {{{}_nC_{r} xy} \over {({}_nC_{n-r} (x \log x)(y \log y))^{-1}}} 87 = ({}_nC_{n-r})^2 \sum_{k= 0}^{\infty}({1 \over x \log x} - {1 \over y \log y})d{1 \over nxy} \times {xy} = \sum_{k=0}^{\infty} a_k f^k = \alpha } これまでのデータベース化された機能のもとでもある方程式たちを構造体として、 まとめて、=> [tuplespace] としてポインタを当てる。 _ struct_ :asperal equation.emerged => [tuplespace] tuplespace.cognitive_system => development -> Omega.Database[import] value.equation_emerged.exclude >- Omega.Database[tuplespace] この連想ポインタは tuplespace 自体をオブジェクト化して、レシーバの.cognitive_system から 実装段階で、Omega データベース化する。 このデータベースの仮の方程式、未定義な式をデータベースから多様体の仕組みを利用して、 データベースから分岐して、value のオブジェクトとしてポインタを当てる。 以下のソースコードは、今まで扱っている多様体のデータを使って、アプリケーションプログラミングとし て、即席スクリプト言語を DSL として書いている。 Omega::DataBase <-> virtual_connect(VIRTUALMACHINE) { blidge_base.network => localmachine.attachment :=> { dhcp.etc_load_file(this.klass) {|list| list.connect[XWin.display _ <- xhost.in(regexpt.pattern)] { ultranetwork.def _struct { asperal_language :this.network_address.included[type.system_pattern] {|regexpt.pattern| <- w.scan |each_string| <= { ipv4.file :file.port subnetmask :file.address file.port <=> file.address FILE *pointer int,char,str :emerge.exclude > array[] BTE.each_string <-> regexpt.pattern { development => file.to_excluded file.scan => regexpt.pattern this.iterator <-> each_string 88 file.reloded => [asperal_language.rebuild] } } } } } } } } class Ultranetwork def virtual_connect load :file => { asperal :virtual_machine.attachment { system.require :file.attachment <- |list.file| :=> { tk.mainloop <- [XWin -multiwindow] startx => file.load.environment in { [blidge_base | host_base].connect(wmware.dhcp) net_work.connect.used[wireshark.demand => exclude(file)] } } } } end def ¡ method として、メソッドを def へリファクタリング機能をつかって、def へと以下の method たちは 取り込まれる。これの作用は、def が one class 並の等号シングレトンとして機能する。 def < blidge_base.network.connect { dhcp.start => { host_name <-> localhost_name {|list| list.exist(connect_type) { <- : tty :xhost -display => list.exist [virtual_connect].list->host :terminal } } } end def < host_base.ethernet.connect { host_name.connect => local_network } } 89 end def < etc.load_file { etc.include(inetd.rc) { virtual_connect(VIRTUAL_MACHINE){|list| list.attachment(etc.load_file) } } } end mainloop{ def.virtual_connect => xhost.localmachine { xhost.client <-> xhost.server } def.network.type <- [Omega.DATABASE] end def.etc.load_file.attachment(VIRTUAL_MACHINE) end end end class UltraNetwork::DATABASE import OMEGA.TUPLESPACE def load_file >- VIRTUAL_MACHINE { in . => attachment_device |for| for.load -> acceptance.hardware virtual_machine.new { tk.loop-> start XWin -multiwindow if dwm <-> new_xwin.start localhost :xhost :display -x xdisplay :-> [preset :XFree.demand>=needed for.set_up install_process >- tar -xvfz "#{load_file}" <-> install_attachment ] else if only :new_xwin.start localhost :xhost :multiwindow . { in display -x attachment :localhost -client from -client into server.XWin -attachment} end condition :{ in .=> check->[xdisplay.install_process]} end 90 def < network_rout wireshark.start -> ethernet.device >- define rout rout.ipstate do |file| file.type <- encoding XWin -filesystem file.included >- make kernel_system.rebuild file.vmware.start do |rout| rout.blidgebase | rout.hostbase -> file.install file.address_ipstate => {"{file}" :=> dwm.state_presense virtual_machine.included[file] } end def < launcher_application network_rout.new |file| file.attachment => { in . new_xwin.start :=> file.included demand.file <- success_exit} end def < terminal_port network_rout.new launcher_application.new |rout| rout.acceptance { vmware.state.process |new_rout| new_rout : attachment.class <-> dwm.state_attachment new_rout -> condition.start_wmware.process} end def < kterm_port launcher_application.new def.included[DATABASE] |rout| rout.attachment <- |new_rout| new_rout.attachment do install.condition < rout.def.terminal_port.exclude[file] end main_loop :file do kterm_port.excluded :=> VIRTUAL_MACHINE |new_rout| start do rout.process -> network_rout.rout [ file,launcher_application, terminal_port, kterm_port].def < included |file| file.all_attachment: file_type :=> encoding-utf8 end end 91 class < def { pholograph_data[] = [R,V,S,E,U,M_n,Z_n,Q,C,N,f,g] source_array <- pholograph_data[] } def > operator_data[] = {nabla,nabla_i nabla_j,Delta,partial, d, int, cap,cup,ni,in,chi,oplus,otimes,bigoplus,bigotimes,d /over df, end def > manifold_emerge c = def.inject >- source_array times def.operator_data[] repository_data <=> c{ c.scan(/tupplespace[]/) import |list| list{ kerf = -2 \int (R + nabla_i nabla_j f)^2e^{-f}dV kerf / imf =< {d \over df}F} } equals_data =~ /list/ list.match(/"#{c}"/) {|list| list.delete jisyo_data_mathmatics <=> list{ list.emerge => {asperal function >- pholograph_data[] times repository_data =< list.update} } ln -s operator_named <= {list} define _struct |list| -> list.element -> manifold_emerge => list.reconstruct > def.inject /^"#{pattern}"/} end import Omega::Tuplespace < Database { {\bigoplus \nabla M^{+}_{-}}.equation_create -> asperal :variable[array] :=> [cognitive_system <-> def < VIRTUALMACHINE.terminal { [ipv4.bloadcast.address : ipv4.network.adress].subnetmask <-> file.port.transport_import : Omega[tuplespace] } } _struct _ Omega[tuplespace] >> VIRTUALMACHINE.terminal.value class < def.VIRTUALMACHINE.system_environment 92 file.reload[hardware] => file.exclude >> file.attachment {=> |file| file.port(wireshark.rout <-> {file.port.transport_export :=> Omega[tuplespace]} assembly_process.file.included >- file.reloaded :- |file.environment| {=> file.type? :=> exist file.regexpt.pattern[scan.flex] => |pattern| <-> file.[scan.compiler] } end end file << } } Omega::Database[tuplespace] { cognitive_system |: -> { DATABASE.create.regexpt_pattern >cognitive_system[tuplespace].recreated >- : =< DATABASE.value >> system_require.application.reloaded[tuplespace] } : _struct _ def.VIRTUALMACHINE.terminal >> { ||machine.attachment|| <-> OBJECT.shift => system.reloaded . in { : _struct _ class.import :-> require mechanics.DATABASE {|regexpt_pattern| :|-> aspective _union _ def _union _} } } end } system.require <- import library.DATABASE { Omega[tuplespace] { cognitive_system : VIRTUALMACHINE.equality_realized {|regexpt_pattern| => value | key [ > cognitive_system.loop.stdout] value : display -bash :xhost -number XWin.terminal key : registry.edit :=> {[cognitive_system.reloaded]} } } } 93 _union _ => DATABASE[tuplespace].aspective_reloaded _union _ :fx | -> |regexpt_pattern| => { VIRTUALMACHIE.recreated-> _union _ | _struct _ def.DATABASE.recreated <- fx >> DATABASE[tuplespace].rebuild } DATABASE[tuplespace] -< {[ > aimed.compiler | aimed.interpreter] | btree.def.distributed >aimed[tuplespace]} aimed[tuplespace] -< btree.class.hyperrout_ struct _ => Omega::Database[tuplespace].value sheap_ union _ :aspective | -> Omega[tuplespace]: | aimed[tuplespace].differented_review } aimed[tuplespace].process => DATABASE[tuplespace].reloaded aimed.different | aimed.stdout >> vale | key [ > cognitive_system.loop.stdin] {|pattern| pattern.scan(value : aimed[def.value] key : aimed[def.key]) } _ struct _ : flex | interpreter.system => expression.iterator[def.first,def.second,def.third,def.fourth] { def < Omega[tuplespace] def.cognitive_system |: -> DATABASE[tuplespace] | aimed[tuplespace] } Omega::Tuplespace < DATABASE {=> norm[Fx] -> . in for def.all_included < aimed[tuplespace].each_scan([regexpt_pattern] <-> DATABASE[tuplespace]) << streem database.excluded >- more_pattern.scan(value : aimed[def.value] key :aimed[def.key]) . in { _struct _ :flex | interpreter.system => expression.iterator[def.all.each -> |value, key| included >- norm[Fx]|[DATABASE[tuplespace] ,aimed[tupespace]] | finality : aimed[tuplespace], DATABASE[tuplespace] : -> def.included(in_all) { def.key | def,value => [DATABASE].recompile & make install : in_all -> _struct _ :aspective :tuplespace : all_homology_created} } } def < Omega::Tuplespace[DATABASE] def.iterator -> |klass,define_method,constant,variable,infinity_data : -> finite_data| 94 def.each_klass?{|value, key| _struct _ :aspective -> tuplespace :all_homology_recreated :make menuconfig {=+ def.key -> aimed[def.key],def.value -> aimed[def.value] {|list| list.developed => <key,value> | <aimed[\n\end{equation}\n\begin{equation}\label{eq:ex7}\n’] -> _union _ :value,key : _struct _ <- (_union _ <-> _struct _ +) begin def.key <-> aimed[value] case :one_ exist :other :bug { result <-> def.key { differented :DATABASE[tuplespace] } return :tuplespace.value.shift -> included<tuplespace> else if :other :bug { success_exit <- bug[value] { cognitive_system.scan(bug[value]) { {[e^{-f}[{2 \int (R + \nablaf^2) \over -(R + \Delta f)}e^{-f}dV} .created_field {=> regexpt.pattern \native_function <-> euler-equation {\n\end{equation}\n\begin{equation}\label{eq:ex8}\n’ all_included <- def.key <-> aimed[value]\n\end{equation}\n\begin{equation}\label{eq:ex9}\nvariable in\n\end{equation}\n\begin{equation}\label{eq:ex10}\nstdout} .developed >= { ping localhost -> blidgebase <-> hostbase.virtualmachine.attachment { xhost :display -> streem_style.value networkconnect.hostbase -> localarea.virtualmachine } :connected -> networkrout : flow_to :localhost.attachment } 96 } _struct : def < hostbase.virtualmachine.attachment => : networkrout.area.build @reviser <-> def.add [ < _struct] @reviser : def.each{listmenu -> listlink | unlinklist > [developed -> {def.key , def.value}.cur @reviser <-> def.rebuild [ < _struct] @reviser.def.<value|key>networkrout-> def.present def.present.flow_to -> hostbase.rout << networkrout.data.<value|key> XWin -multiwindow <-> networkrout.data[\n\end{equation}\n\begin{equation}\label{eq:ex11}\n’] def <\n\end{equation}\n\begin{equation}\label{eq:ex12}\n‘:\n\end{equation}\n\section{Derived hypothesis equations}\n\begin{equation}\n. \bigoplus \nabla C^{+}_{-} \cong M_3, R \supset Q, R \cap Q, R \subset M_3, C^{+} \bigoplus M_n, E^{+} \cap R^{+} \quad (extracted)\\ G(s) := \int_M K(x,s) \Gamma(\alpha(x,s)) \; d\mu(x)\n\end{equation}\n\begin{equation}\nH(x) := (1/N)\sum_{i=1}^N \phi_i(x) + \prod_{i=1}^N \psi_i(x)  \quad (Higgs-like ansatz)\n\end{equation}\n\begin{equation}\nD_s G(s) + \lambda \int_M H(x) K(x,s) d\mu(x) + R_s = 0  (complexity=3)\n\end{equation}\n\begin{equation}\n \quad (extracted)\\ G(s) := \int_M K(x,s) \Gamma(\alpha(x,s)) \; d\mu(x)\n\end{equation}\n\begin{equation}\nH(x) := (1/N)\sum_{i=1}^N \phi_i(x) + \prod_{i=1}^N \psi_i(x)  \quad (Higgs-like ansatz)\n\end{equation}\n\begin{equation}\nD_s G(s) + \lambda \int_M H(x) K(x,s) d\mu(x) + R_s = 0  (complexity=3)\n\end{equation}\n\begin{equation}\n. {\nabla \over \Delta} \int x f(x) dx, {\nabla R \over \Delta f}, \square = 2{\int {(R + \nabla_{i} \nabla_{j} f)^2 \over -(R + \Delta f)}}e^{-f}dV \square = {\nabla R \over \Delta f}, {d \over dt}g_{ij} = \square \to {\nabla f \over \Delta x}, (R + |\nabla f|^2)dm \to -2(R + \nabla_{i} \nabla_{j} f)^2 e^{-f}dV x^n + y^n = z^n \to \nabla \psi^2 = 8 \pi G T^{\mu\nu}, f(x + y) \ge f(x) \circ f(y) \mathrm{im}f / \mathrm{ker}f = \partial f, \mathrm{ker}f = \partial f, \mathrm{ker}f / \mathrm{im}f \cong \partial f, \mathrm{ker}f = f^{-1}(x)xf(x) f^{-1}(x)xf(x) = \int \partial f(x) d(\mathrm{kerf}) \to \nabla f = 2 _{n}C_{r} = {}_{n}C_{n-r} \to \mathrm{im}f / \mathrm{ker}f \cong \mathrm{ker}f / \mathrm{im}f \quad (extracted)\\ G(s) := \int_M K(x,s) \Gamma(\alpha(x,s)) \; d\mu(x)\n\end{equation}\n\begin{equation}\nH(x) := (1/N)\sum_{i=1}^N \phi_i(x) + \prod_{i=1}^N \psi_i(x)  \quad (Higgs-like ansatz)\n\end{equation}\n\begin{equation}\nD_s G(s) + \lambda \int_M H(x) K(x,s) d\mu(x) + R_s = 0  (complexity=3)\n\end{equation}\n\begin{equation}\n. this equation \quad (extracted)\\ G(s) := \int_M K(x,s) \Gamma(\alpha(x,s)) \; d\mu(x)\n\end{equation}\n\begin{equation}\nH(x) := (1/N)\sum_{i=1}^N \phi_i(x) + \prod_{i=1}^N \psi_i(x)  \quad (Higgs-like ansatz)\n\end{equation}\n\begin{equation}\nD_s G(s) + \lambda \int_M H(x) K(x,s) d\mu(x) + R_s = 0  (complexity=3)\n\end{equation}\n\begin{equation}\n. V/W = R/C \sum^{\infty}_{k=0}{x^k \over a_k f^k}, W/V = C/R \sum^{\infty}_{k=0}{a_k f^k \over x^k} V/W \cong W/V \cong R/C(\sum^{\infty}_{r=0} {}_nC_{r})^{-1} \sum^{\infty}_{k=0} x^k This equation is diffrential equation, then \quad (extracted)\\ G(s) := \int_M K(x,s) \Gamma(\alpha(x,s)) \; d\mu(x)\n\end{equation}\n\begin{equation}\nH(x) := (1/N)\sum_{i=1}^N \phi_i(x) + \prod_{i=1}^N \psi_i(x)  \quad (Higgs-like ansatz)\n\end{equation}\n\begin{equation}\nD_s G(s) + \lambda \int_M H(x) K(x,s) d\mu(x) + R_s = 0  (complexity=3)\n\end{equation}\n\begin{equation}\nis included with \quad (extracted)\\ G(s) := \int_M K(x,s) \Gamma(\alpha(x,s)) \; d\mu(x)\n\end{equation}\n\begin{equation}\nH(x) := (1/N)\sum_{i=1}^N \phi_i(x) + \prod_{i=1}^N \psi_i(x)  \quad (Higgs-like ansatz)\n\end{equation}\n\begin{equation}\nD_s G(s) + \lambda \int_M H(x) K(x,s) d\mu(x) + R_s = 0  (complexity=3)\n\end{equation}\n\begin{equation}\nW/V = xF(x), \chi(x) = (-1)^k a_k, \\Gamma(x) = \int e^{-x} x^{1 -t}dx, \sum^{n}_{k=0}a_k f^k = (f^k)’ \sum^{n}_{k=0}a_k f^k = \sum^{\infty}_{k~0} {}_{n}C_{r} f^k = (f^k)’, \sum^{\infty}_{k=0} a_k f^k = [f(x)], \sum^{\infty}_{k=0} a_k f^k = \alpha, \sum^{\infty}_{k=0} {1 \over a_k f^k}, \sum^{\infty}_{k=0} (a_k f^k)^{-1} = {1 \over 1 - z} {\int \int {1 \over (x \log x)(y \log y)}dxy} = {{{}_nC_{r} xy} \over {({}_nC_{n-r} (x \log x)(y \log y))^{-1}}} 87 = ({}_nC_{n-r})^2 \sum_{k= 0}^{\infty}({1 \over x \log x} - {1 \over y \log y})d{1 \over nxy} \times {xy} = \sum_{k=0}^{\infty} a_k f^k = \alpha } これまでのデータベース化された機能のもとでもある方程式たちを構造体として、 まとめて、=> [tuplespace] としてポインタを当てる。 _ struct_ :asperal equation.emerged => [tuplespace] tuplespace.cognitive_system => development -> Omega.Database[import] value.equation_emerged.exclude >- Omega.Database[tuplespace] この連想ポインタは tuplespace 自体をオブジェクト化して、レシーバの.cognitive_system から 実装段階で、Omega データベース化する。 このデータベースの仮の方程式、未定義な式をデータベースから多様体の仕組みを利用して、 データベースから分岐して、value のオブジェクトとしてポインタを当てる。 以下のソースコードは、今まで扱っている多様体のデータを使って、アプリケーションプログラミングとし て、即席スクリプト言語を DSL として書いている。 Omega::DataBase <-> virtual_connect(VIRTUALMACHINE) { blidge_base.network => localmachine.attachment :=> { dhcp.etc_load_file(this.klass) {|list| list.connect[XWin.display _ <- xhost.in(regexpt.pattern)] { ultranetwork.def _struct { asperal_language :this.network_address.included[type.system_pattern] {|regexpt.pattern| <- w.scan |each_string| <= { ipv4.file :file.port subnetmask :file.address file.port <=> file.address FILE *pointer int,char,str :emerge.exclude > array[] BTE.each_string <-> regexpt.pattern { development => file.to_excluded file.scan => regexpt.pattern this.iterator <-> each_string 88 file.reloded => [asperal_language.rebuild] } } } } } } } } class Ultranetwork def virtual_connect load :file => { asperal :virtual_machine.attachment { system.require :file.attachment <- |list.file| :=> { tk.mainloop <- [XWin -multiwindow] startx => file.load.environment in { [blidge_base | host_base].connect(wmware.dhcp) net_work.connect.used[wireshark.demand => exclude(file)] } } } } end def ¡ method として、メソッドを def へリファクタリング機能をつかって、def へと以下の method たちは 取り込まれる。これの作用は、def が one class 並の等号シングレトンとして機能する。 def < blidge_base.network.connect { dhcp.start => { host_name <-> localhost_name {|list| list.exist(connect_type) { <- : tty :xhost -display => list.exist [virtual_connect].list->host :terminal } } } end def < host_base.ethernet.connect { host_name.connect => local_network } } 89 end def < etc.load_file { etc.include(inetd.rc) { virtual_connect(VIRTUAL_MACHINE){|list| list.attachment(etc.load_file) } } } end mainloop{ def.virtual_connect => xhost.localmachine { xhost.client <-> xhost.server } def.network.type <- [Omega.DATABASE] end def.etc.load_file.attachment(VIRTUAL_MACHINE) end end end class UltraNetwork::DATABASE import OMEGA.TUPLESPACE def load_file >- VIRTUAL_MACHINE { in . => attachment_device |for| for.load -> acceptance.hardware virtual_machine.new { tk.loop-> start XWin -multiwindow if dwm <-> new_xwin.start localhost :xhost :display -x xdisplay :-> [preset :XFree.demand>=needed for.set_up install_process >- tar -xvfz "#{load_file}" <-> install_attachment ] else if only :new_xwin.start localhost :xhost :multiwindow . { in display -x attachment :localhost -client from -client into server.XWin -attachment} end condition :{ in .=> check->[xdisplay.install_process]} end 90 def < network_rout wireshark.start -> ethernet.device >- define rout rout.ipstate do |file| file.type <- encoding XWin -filesystem file.included >- make kernel_system.rebuild file.vmware.start do |rout| rout.blidgebase | rout.hostbase -> file.install file.address_ipstate => {"{file}" :=> dwm.state_presense virtual_machine.included[file] } end def < launcher_application network_rout.new |file| file.attachment => { in . new_xwin.start :=> file.included demand.file <- success_exit} end def < terminal_port network_rout.new launcher_application.new |rout| rout.acceptance { vmware.state.process |new_rout| new_rout : attachment.class <-> dwm.state_attachment new_rout -> condition.start_wmware.process} end def < kterm_port launcher_application.new def.included[DATABASE] |rout| rout.attachment <- |new_rout| new_rout.attachment do install.condition < rout.def.terminal_port.exclude[file] end main_loop :file do kterm_port.excluded :=> VIRTUAL_MACHINE |new_rout| start do rout.process -> network_rout.rout [ file,launcher_application, terminal_port, kterm_port].def < included |file| file.all_attachment: file_type :=> encoding-utf8 end end 91 class < def { pholograph_data[] = [R,V,S,E,U,M_n,Z_n,Q,C,N,f,g] source_array <- pholograph_data[] } def > operator_data[] = {nabla,nabla_i nabla_j,Delta,partial, d, int, cap,cup,ni,in,chi,oplus,otimes,bigoplus,bigotimes,d /over df, end def > manifold_emerge c = def.inject >- source_array times def.operator_data[] repository_data <=> c{ c.scan(/tupplespace[]/) import |list| list{ kerf = -2 \int (R + nabla_i nabla_j f)^2e^{-f}dV kerf / imf =< {d \over df}F} } equals_data =~ /list/ list.match(/"#{c}"/) {|list| list.delete jisyo_data_mathmatics <=> list{ list.emerge => {asperal function >- pholograph_data[] times repository_data =< list.update} } ln -s operator_named <= {list} define _struct |list| -> list.element -> manifold_emerge => list.reconstruct > def.inject /^"#{pattern}"/} end import Omega::Tuplespace < Database { {\bigoplus \nabla M^{+}_{-}}.equation_create -> asperal :variable[array] :=> [cognitive_system <-> def < VIRTUALMACHINE.terminal { [ipv4.bloadcast.address : ipv4.network.adress].subnetmask <-> file.port.transport_import : Omega[tuplespace] } } _struct _ Omega[tuplespace] >> VIRTUALMACHINE.terminal.value class < def.VIRTUALMACHINE.system_environment 92 file.reload[hardware] => file.exclude >> file.attachment {=> |file| file.port(wireshark.rout <-> {file.port.transport_export :=> Omega[tuplespace]} assembly_process.file.included >- file.reloaded :- |file.environment| {=> file.type? :=> exist file.regexpt.pattern[scan.flex] => |pattern| <-> file.[scan.compiler] } end end file << } } Omega::Database[tuplespace] { cognitive_system |: -> { DATABASE.create.regexpt_pattern >cognitive_system[tuplespace].recreated >- : =< DATABASE.value >> system_require.application.reloaded[tuplespace] } : _struct _ def.VIRTUALMACHINE.terminal >> { ||machine.attachment|| <-> OBJECT.shift => system.reloaded . in { : _struct _ class.import :-> require mechanics.DATABASE {|regexpt_pattern| :|-> aspective _union _ def _union _} } } end } system.require <- import library.DATABASE { Omega[tuplespace] { cognitive_system : VIRTUALMACHINE.equality_realized {|regexpt_pattern| => value | key [ > cognitive_system.loop.stdout] value : display -bash :xhost -number XWin.terminal key : registry.edit :=> {[cognitive_system.reloaded]} } } } 93 _union _ => DATABASE[tuplespace].aspective_reloaded _union _ :fx | -> |regexpt_pattern| => { VIRTUALMACHIE.recreated-> _union _ | _struct _ def.DATABASE.recreated <- fx >> DATABASE[tuplespace].rebuild } DATABASE[tuplespace] -< {[ > aimed.compiler | aimed.interpreter] | btree.def.distributed >aimed[tuplespace]} aimed[tuplespace] -< btree.class.hyperrout_ struct _ => Omega::Database[tuplespace].value sheap_ union _ :aspective | -> Omega[tuplespace]: | aimed[tuplespace].differented_review } aimed[tuplespace].process => DATABASE[tuplespace].reloaded aimed.different | aimed.stdout >> vale | key [ > cognitive_system.loop.stdin] {|pattern| pattern.scan(value : aimed[def.value] key : aimed[def.key]) } _ struct _ : flex | interpreter.system => expression.iterator[def.first,def.second,def.third,def.fourth] { def < Omega[tuplespace] def.cognitive_system |: -> DATABASE[tuplespace] | aimed[tuplespace] } Omega::Tuplespace < DATABASE {=> norm[Fx] -> . in for def.all_included < aimed[tuplespace].each_scan([regexpt_pattern] <-> DATABASE[tuplespace]) << streem database.excluded >- more_pattern.scan(value : aimed[def.value] key :aimed[def.key]) . in { _struct _ :flex | interpreter.system => expression.iterator[def.all.each -> |value, key| included >- norm[Fx]|[DATABASE[tuplespace] ,aimed[tupespace]] | finality : aimed[tuplespace], DATABASE[tuplespace] : -> def.included(in_all) { def.key | def,value => [DATABASE].recompile & make install : in_all -> _struct _ :aspective :tuplespace : all_homology_created} } } def < Omega::Tuplespace[DATABASE] def.iterator -> |klass,define_method,constant,variable,infinity_data : -> finite_data| 94 def.each_klass?{|value, key| _struct _ :aspective -> tuplespace :all_homology_recreated :make menuconfig {=+ def.key -> aimed[def.key],def.value -> aimed[def.value] {|list| list.developed => <key,value> | <aimed[ \quad (extracted)\\ G(s) := \int_M K(x,s) \Gamma(\alpha(x,s)) \; d\mu(x)\n\end{equation}\n\begin{equation}\nH(x) := (1/N)\sum_{i=1}^N \phi_i(x) + \prod_{i=1}^N \psi_i(x)  \quad (Higgs-like ansatz)\n\end{equation}\n\begin{equation}\nD_s G(s) + \lambda \int_M H(x) K(x,s) d\mu(x) + R_s = 0  (complexity=3)\n\end{equation}\n\begin{equation}\n’] -> _union _ :value,key : _struct _ <- (_union _ <-> _struct _ +) begin def.key <-> aimed[value] case :one_ exist :other :bug { result <-> def.key { differented :DATABASE[tuplespace] } return :tuplespace.value.shift -> included<tuplespace> else if :other :bug { success_exit <- bug[value] { cognitive_system.scan(bug[value]) { {[e^{-f}[{2 \int (R + \nablaf^2) \over -(R + \Delta f)}e^{-f}dV} .created_field {=> regexpt.pattern \native_function <-> euler-equation { \quad (extracted)\\ G(s) := \int_M K(x,s) \Gamma(\alpha(x,s)) \; d\mu(x)\n\end{equation}\n\begin{equation}\nH(x) := (1/N)\sum_{i=1}^N \phi_i(x) + \prod_{i=1}^N \psi_i(x)  \quad (Higgs-like ansatz)\n\end{equation}\n\begin{equation}\nD_s G(s) + \lambda \int_M H(x) K(x,s) d\mu(x) + R_s = 0  (complexity=3)\n\end{equation}\n\begin{equation}\n’ all_included <- def.key <-> aimed[value] \quad (extracted)\\ G(s) := \int_M K(x,s) \Gamma(\alpha(x,s)) \; d\mu(x)\n\end{equation}\n\begin{equation}\nH(x) := (1/N)\sum_{i=1}^N \phi_i(x) + \prod_{i=1}^N \psi_i(x)  \quad (Higgs-like ansatz)\n\end{equation}\n\begin{equation}\nD_s G(s) + \lambda \int_M H(x) K(x,s) d\mu(x) + R_s = 0  (complexity=3)\n\end{equation}\n\begin{equation}\nvariable in \quad (extracted)\\ G(s) := \int_M K(x,s) \Gamma(\alpha(x,s)) \; d\mu(x)\n\end{equation}\n\begin{equation}\nH(x) := (1/N)\sum_{i=1}^N \phi_i(x) + \prod_{i=1}^N \psi_i(x)  \quad (Higgs-like ansatz)\n\end{equation}\n\begin{equation}\nD_s G(s) + \lambda \int_M H(x) K(x,s) d\mu(x) + R_s = 0  (complexity=3)\n\end{equation}\n\begin{equation}\nstdout} .developed >= { ping localhost -> blidgebase <-> hostbase.virtualmachine.attachment { xhost :display -> streem_style.value networkconnect.hostbase -> localarea.virtualmachine } :connected -> networkrout : flow_to :localhost.attachment } 96 } _struct : def < hostbase.virtualmachine.attachment => : networkrout.area.build @reviser <-> def.add [ < _struct] @reviser : def.each{listmenu -> listlink | unlinklist > [developed -> {def.key , def.value}.cur @reviser <-> def.rebuild [ < _struct] @reviser.def.<value|key>networkrout-> def.present def.present.flow_to -> hostbase.rout << networkrout.data.<value|key> XWin -multiwindow <-> networkrout.data[ \quad (extracted)\\ G(s) := \int_M K(x,s) \Gamma(\alpha(x,s)) \; d\mu(x)\n\end{equation}\n\begin{equation}\nH(x) := (1/N)\sum_{i=1}^N \phi_i(x) + \prod_{i=1}^N \psi_i(x)  \quad (Higgs-like ansatz)\n\end{equation}\n\begin{equation}\nD_s G(s) + \lambda \int_M H(x) K(x,s) d\mu(x) + R_s = 0  (complexity=3)\n\end{equation}\n\begin{equation}\n’] def < \quad (extracted)\\ G(s) := \int_M K(x,s) \Gamma(\alpha(x,s)) \; d\mu(x)\n\end{equation}\n\begin{equation}\nH(x) := (1/N)\sum_{i=1}^N \phi_i(x) + \prod_{i=1}^N \psi_i(x)  \quad (Higgs-like ansatz)\n\end{equation}\n\begin{equation}\nD_s G(s) + \lambda \int_M H(x) K(x,s) d\mu(x) + R_s = 0  (complexity=3)\n\end{equation}\n\begin{equation}\n‘: \quad (extracted)\\ G(s) := \int_M K(x,s) \Gamma(\alpha(x,s)) \; d\mu(x)\n\end{equation}\n\begin{equation}\nH(x) := (1/N)\sum_{i=1}^N \phi_i(x) + \prod_{i=1}^N \psi_i(x)  \quad (Higgs-like ansatz)\n\end{equation}\n\begin{equation}\nD_s G(s) + \lambda \int_M H(x) K(x,s) d\mu(x) + R_s = 0  (complexity=3)\n\end{equation}\n\section{Validation plan}\n1) Specify M and measure; discretize.\n2) Propose parametric forms; fit to data/synthetic tests.\n3) Numerical convergence and residual analysis.\n\section{Caution}\nThis is EXPLORATORY/HYPOTHESIS output. Not a proof.\n\end{document}\n
